\chapter{PLANTEAMIENTO DEL EMPRENDIMIENTO PRODUCTIVO}
\label{chap:planteamiento_emprendimiento}

\section{Diagnóstico del contexto productivo}
\label{sec:diagnostico_contexto}
Análisis del entorno social, económico y sectorial en el cual se insertará el emprendimiento productivo, identificando brechas de oferta y necesidades no atendidas.

\section{Objetivos del emprendimiento productivo}
\label{sec:objetivos_emprendimiento}

\subsection{Objetivo general}
Establecer un modelo de negocio sostenible para la producción y comercialización del bien o servicio en el mercado regional.

\subsection{Objetivos específicos}
\begin{itemize}
    \item Realizar el estudio de mercado para identificar la demanda insatisfecha y caracterizar al público objetivo.
    \item Diseñar la estructura organizacional y los procesos operativos de producción/servicio.
    \item Determinar los costos de inversión, capital de operación y precios de venta garantizando rentabilidad.
\end{itemize}

\subsection{Misión (opcional)}
Proveer productos/servicios de alta calidad que satisfagan las necesidades del mercado contribuyendo al desarrollo local.

\subsection{Visión (opcional)}
Consolidarse como una empresa referente en el sector productivo por su innovación, responsabilidad social y excelencia operativa.

\section{Justificación}
\label{sec:justificacion_emprendimiento}
Fundamentación de la importancia económica, social y técnica del emprendimiento, destacando la generación de empleo y aprovechamiento de recursos locales.
